\ifdefined\iscombined
	\lecturetitle{Petri-nets and CTL}
\else
\documentclass{article}
\usepackage{changepage}
\usepackage{centernot}
\usepackage{listings}
\usepackage{amsthm}
\usepackage{braket}
\usepackage{comment}
\usepackage{algorithm}
\usepackage{algpseudocode}
\usepackage{amsmath}
\usepackage{amsfonts}
\usepackage{sectsty}
\usepackage{hyperref}
\usepackage{fancyhdr}
\usepackage[top=1in,bottom=1in,left=1in,right=1in]{geometry}
\usepackage{float}
\usepackage{graphicx}
\usepackage{tabularx}
\usepackage{mathtools}
\usepackage{tikz}
\usetikzlibrary{backgrounds, calc, positioning}

\date{
	\vspace{-.75em}
	{\large \today}
}

\title{
	\vspace*{-1.5em}
	{\Huge Lecture 20} \\
	\vspace{-1em}
	\rule{\textwidth/4}{.01em} \\
	\vspace{-.25em}
	{\Large Petri-nets and CTL}
	\vspace{-.75em}
}

\author{
	{\large Matthew Farah}
}

\hypersetup{
	colorlinks=true,
	linkcolor=blue,
	filecolor=magenta,
	urlcolor=blue,
	citecolor=blue,
	anchorcolor=blue
}

\fancyhead{}
\fancyhead[L]{\today}
\fancyhead[R]{Lecture 20 - Petri-nets and CTL}

\setlength{\parindent}{0cm}

\lstset{
	basicstyle=\ttfamily\small,
	frame=single,
	tabsize=4,
	breaklines=true,
	numbers=left,
	numberstyle=\tiny,
	numbersep=8pt,
	xleftmargin=1.5em,
	framexleftmargin=1.5em
}

\newcounter{example}
\renewcommand{\theexample}{\arabic{example}}

\newenvironment{example}[1][]{
	\refstepcounter{example}
	\begin{adjustwidth}{1.5em}{}
	\noindent\textbf{\ifx&#1&Example \theexample:\else Example \theexample \ (#1):\fi}
	\ignorespaces
}{
	\vspace{.5em}
	\end{adjustwidth}
}

\newcounter{subexample}
\renewcommand{\thesubexample}{\alph{subexample}}

\newenvironment{subexample}{
	\setcounter{step}{0}
	\refstepcounter{subexample}
	\vspace{1em}\par\noindent
	\begin{adjustwidth}{1.5em}{}
	\noindent\textbf{(\thesubexample)}
	\ignorespaces
}{
	\par
	\vspace{1em}
	\end{adjustwidth}
}

\newenvironment{solution}{
	\vspace{.5em}
	\par
	\begin{adjustwidth}{1.5em}{}
	\textbf{{Solution:}}
	\par
	\vspace{.5em}
}{
	\vspace{1em}
	\end{adjustwidth}
}

\newcounter{step}
\renewcommand{\thestep}{\Roman{step}}

\newenvironment{step}{
	\refstepcounter{step}
	\vspace{1em}\par\noindent
	\ignorespaces\noindent\textbf{(\thestep)}
	\ignorespaces
}{
	\par
	\vspace{1em}
}

\sectionfont{\fontsize{12}{12}\bfseries}
\subsectionfont{\fontsize{10}{12}\normalfont}
\subsubsectionfont{\fontsize{10}{12}\normalfont}

\begin{document}

\maketitle
\pagestyle{fancy}

\fi

\begin{itemize}
	\item Petri-nets are used to model non-determinism in software.
	\item The critical section of a program refers to where shared resources are stored.
	\begin{itemize}
		\item Semaphores are used to control processes' access to the critical section.
	\end{itemize}
	\item Inhibator arcs are arcs (arrows) that can only be taken if its input place is empty.
	\item Transitions may have different priorities.
	\begin{itemize}
		\item Helps to solve conflicts and prevent deadlocks.
	\end{itemize}
	\item The start state of a Petri-net describes its initial arrangement of tokens before any transitions are taken.
	\item The state of a Petri-net describes an valid arrangement of its tokens reachable from the start-state
	\item A path through a Petri-net is a valid sequence of its states.
	\item Computational Tree Language (CTL) is a query language.
	\begin{itemize}
		\item Frequently used for evaluating properties of Petri-nets.
		\item Evaluate to true or false.
	\end{itemize}
	\item CTL statements are use two types of quantifiers:
	\begin{enumerate}
		\item General quantifiers.
		\begin{itemize}
			\item Quantifies over all paths from the current state.
		\end{itemize}
		\item Path-specific quantifiers.
		\begin{itemize}
			\item Quantifies along a given path from the current state.
		\end{itemize}
	\end{enumerate}
	\item There are two general quantifiers in CTL:
	\begin{enumerate}
		\item $A$.
		\begin{itemize}
			\item Means ``for all paths in the model''.
			\item \emph{Aside: Think of it like $\forall$}.
		\end{itemize}
		\item $E$.
		\begin{itemize}
			\item Means ``there exists a path in the mode''.
			\item \emph{Aside: Think of it like $\exists$}.
		\end{itemize}
	\end{enumerate}
	\item There are four path-specific quantifiers in CTL:
	\begin{enumerate}
		\item $F$.
		\begin{itemize}
			\item Means ``finally''.
			\item \emph{Aside: Think of it like the property should eventually hold}.
		\end{itemize}
		\item $G$.
		\begin{itemize}
			\item Means ``globally''.
			\item \emph{Aside: Think of it as meaning the property should hold for all valid states}.
		\end{itemize}
		\item $X$.
		\begin{itemize}
			\item Means ``next''.
		\end{itemize}
		\item $U$.
		\begin{itemize}
			\item Means ``until''.
			\item \emph{Aside: Think of it as meaning the property should hold \underline{at least} until some other property is true (the property can still hold after the other becomes true)}.
		\end{itemize}
	\end{enumerate}
	\item The criteria for these quantifiers to be satisfied looks like the following when graphed, where:
	\begin{itemize}
		\item $p$ and $q$ are properties of the system being queries for.
		\item Blue vertices correspond to states where the query for $p$ is true.
		\item Red vertices correspond to states where the query for $q$ is true.
		\item Black vertices correspond to states where the query may or not be true (doesn't matter which).
	\end{itemize}
\end{itemize}

\vspace{1em}

\definecolor{bp}{RGB}{50,50,220}      % Blue
\definecolor{rp}{RGB}{220,50,70}      % Red

\tikzset{
	nd/.style  = {circle, fill=black, inner sep=0pt, minimum size=4.5pt},
	bn/.style  = {circle, fill=bp,    inner sep=0pt, minimum size=5.5pt},
	rn/.style  = {circle, fill=rp,    inner sep=0pt, minimum size=5.5pt},
	eg/.style  = {thick, black!55},
	be/.style  = {very thick, bp},
	re/.style  = {very thick, rp},
}

\newcommand{\treecoords}{
	\coordinate (N0)  at ( 0.0,  0.0);
	\coordinate (N1)  at (-1.2, -1.0);
	\coordinate (N2)  at ( 1.2, -1.0);
	\coordinate (N3)  at (-1.8, -2.0);
	\coordinate (N4)  at (-0.6, -2.0);
	\coordinate (N5)  at ( 0.6, -2.0);
	\coordinate (N6)  at ( 1.8, -2.0);
	\coordinate (N7)  at (-2.1, -3.0);
	\coordinate (N8)  at (-1.5, -3.0);
	\coordinate (N9)  at (-0.9, -3.0);
	\coordinate (N10) at (-0.3, -3.0);
	\coordinate (N11) at ( 0.3, -3.0);
	\coordinate (N12) at ( 0.9, -3.0);
	\coordinate (N13) at ( 1.5, -3.0);
	\coordinate (N14) at ( 2.1, -3.0);
}

\newcommand{\treebg}{
	(0, 0.30) -- (-2.45, -3.25) -- (2.45, -3.25) -- cycle;
}

\newcommand{\treeedges}{
	\foreach \p/\c in
		{0/1,0/2,1/3,1/4,2/5,2/6,3/7,3/8,4/9,4/10,5/11,5/12,6/13,6/14}{
		\draw[eg] (N\p) -- (N\c);
	}
}

\newcommand{\allblack}{
	\foreach \i in {0,...,14}{ \node[nd] at (N\i) {}; }
}

%NOTE: Most of this was generated by Claude.
% I can confirm it is correct however.
% It's based on graphic given in lecture slides and I edited the .tikz substantially it to fix all the mistakes it made.
% AGI my ass.

\begin{center}
	\scalebox{.65}{
		\begin{tikzpicture}

		% =========== grid constants ===========
		\pgfmathsetmacro{\ColSep}{5.8}
		\pgfmathsetmacro{\RowSep}{5.6}

		% =========== column headers ===========
		\node[font=\large\sffamily\bfseries, anchor=south] at (0,1.0)
			{F\,\textcolor{bp}{p}};
		\node[font=\small\sffamily, anchor=south] at (0,0.5)
			{finally \textcolor{bp}{p}};

		\node[font=\large\sffamily\bfseries, anchor=south] at (\ColSep,1.0)
			{G\,\textcolor{bp}{p}};
		\node[font=\small\sffamily, anchor=south] at (\ColSep,0.5)
			{globally \textcolor{bp}{p}};

		\node[font=\large\sffamily\bfseries, anchor=south] at (2*\ColSep,1.0)
			{X\,\textcolor{bp}{p}};
		\node[font=\small\sffamily, anchor=south] at (2*\ColSep,0.5)
			{next \textcolor{bp}{p}};

		\node[font=\large\sffamily\bfseries, anchor=south] at (3*\ColSep,1.0)
			{\textcolor{bp}{p}\;U\;\textcolor{rp}{q}};

		\node[font=\small\sffamily, anchor=south] at (3*\ColSep,0.5)
			{\textcolor{bp}{p} until \textcolor{rp}{q}};

		% =========== row headers ===========
		\node[font=\LARGE\sffamily\bfseries, anchor=east] at (-2.9, -1.5) {A};
		\node[font=\LARGE\sffamily\bfseries, anchor=east] at (-2.9, {-1.5-\RowSep}) {E};

		% ================================================================
		%  ROW 1 — A  (for All paths)
		% ================================================================

		% ---- AF p ----
		%  All paths eventually reach p.
		%  p (blue) holds at levels 2-3; every path passes through level 2.
		\begin{scope}[shift={(0,0)}]
			\treecoords \treebg \treeedges

			\foreach \i in {0,1,2}{ \node[nd] at (N\i) {}; }
			\foreach \i in {3,...,6}{ \node[bn] at (N\i) {}; }
			\foreach \i in {7,...,14}{ \node[nd] at (N\i) {}; }
		\end{scope}

		% ---- AG p ----
		%  All paths, p holds globally (everywhere).
		\begin{scope}[shift={(\ColSep,0)}]
			\treecoords \treebg \treeedges
			\foreach \i in {0,...,14}{ \node[bn] at (N\i) {}; }
		\end{scope}

		% ---- AX p ----
		%  All paths, p holds at the next state.
		%  Only the two level-1 children are blue.
		\begin{scope}[shift={(2*\ColSep,0)}]
			\treecoords \treebg \treeedges
			\allblack
			\foreach \i in {1,2}{ \node[bn] at (N\i) {}; }
		\end{scope}

		% ---- A[p U q] ----
		%  All paths: p (blue) holds until q (red) becomes true.
		%  p at root + level 1; q at levels 2-3.
		\begin{scope}[shift={(3*\ColSep,0)}]
			\treecoords \treebg \treeedges
			% blue (p): root, level 1
			\foreach \i in {0,1,2}{ \node[bn] at (N\i) {}; }
			% red (q): levels 2-3
			\foreach \i in {3,...,6}{ \node[rn] at (N\i) {}; }
			\foreach \i in {7,...,14}{ \node[nd] at (N\i) {}; }
		\end{scope}

		% ================================================================
		%  ROW 2 — E  (there Exists a path)
		% ================================================================

		% ---- EF p ----
		%  There exists a path on which p eventually holds.
		%  One node is blue; path root -> T1 -> T4 reaches it.
		\begin{scope}[shift={(0,-\RowSep)}]
			\treecoords \treebg \treeedges
			\allblack
			% highlighted path
			\draw[be] (N0) -- (N1);
			\draw[be] (N1) -- (N4);
			% single blue node at the end
			\node[bn] at (N4) {};
		\end{scope}

		% ---- EG p ----
		%  There exists a path on which p holds globally.
		%  Path root -> T2 -> T5 -> T12, all nodes blue.
		\begin{scope}[shift={(\ColSep,-\RowSep)}]
			\treecoords \treebg \treeedges
			\allblack
			% highlighted edges
			\draw[be] (N0) -- (N2);
			\draw[be] (N2) -- (N5);
			\draw[be] (N5) -- (N12);
			% blue nodes along the path
			\foreach \i in {0,2,5,12}{ \node[bn] at (N\i) {}; }
		\end{scope}

		% ---- EX p ----
		%  There exists a path where the next state satisfies p.
		%  One level-1 child is blue.
		\begin{scope}[shift={(2*\ColSep,-\RowSep)}]
			\treecoords \treebg \treeedges
			\allblack
			% highlighted edge
			\draw[be] (N0) -- (N2);
			% single blue node
			\node[bn] at (N2) {};
		\end{scope}

		% ---- E[p U q] ----
		%  There exists a path where p holds until q.
		%  Path root -> T2 -> T6 -> T14.  Blue (p): root, T2, T6.  Red (q): T14.
		\begin{scope}[shift={(3*\ColSep,-\RowSep)}]
			\treecoords \treebg \treeedges
			\allblack
			% highlighted edges
			\draw[be] (N0) -- (N2);
			\draw[be] (N2) -- (N6);
			\draw[re] (N6) -- (N14);
			% blue (p) nodes
			\foreach \i in {0,2,6}{ \node[bn] at (N\i) {}; }
			% red (q) node
			\node[rn] at (N14) {};
		\end{scope}

		\end{tikzpicture}
	}
\end{center}

\vspace{1em}

\begin{itemize}
	\item The state of a Petri-net can be described by a vector $v$.
	\begin{itemize}
		\item An index $i$ of $v$ represents a place in the Petri-net.
		\item $v[i]$ represents the number of tokens in place $i$ in the given state.
	\end{itemize}
\end{itemize}

\ifdefined\iscombined
	\newpage
\else
	\end{document}
\fi
